\documentclass[11pt]{article}

\usepackage[T1]{fontenc}
\usepackage{lmodern}
\usepackage{amsmath,amssymb,amsthm,mathtools}
\usepackage{microtype}
\usepackage{mathrsfs}
\usepackage[margin=1.05in]{geometry}
\usepackage{enumitem}
\usepackage{booktabs}
\usepackage{xurl}
\usepackage{needspace}
\usepackage[colorlinks=true,linkcolor=black,citecolor=black,urlcolor=blue]{hyperref}
\hypersetup{pdftitle={Microscopic-Macroscopic Independence and Poisson Peeling for Luce Permutations},pdfauthor={Christopher D. Long}}

\numberwithin{equation}{section}
\allowdisplaybreaks
\setlength{\emergencystretch}{2em}

\newcommand{\E}{\mathbb E}
\newcommand{\Pp}{\mathbb P}
\newcommand{\Var}{\operatorname{Var}}
\newcommand{\Pois}{\operatorname{Pois}}
\newcommand{\PD}{\operatorname{PD}}
\newcommand{\GEM}{\operatorname{GEM}}
\newcommand{\TV}{d_{\mathrm{TV}}}
\newcommand{\one}{\mathbf 1}
\newcommand{\dd}{\,\mathrm d}
\newcommand{\cL}{\mathcal L}
\newcommand{\cP}{\mathcal P}
\newcommand{\cC}{\mathcal C}
\newcommand{\eps}{\varepsilon}

\newtheorem{theorem}{Theorem}[section]
\newtheorem{proposition}[theorem]{Proposition}
\newtheorem{lemma}[theorem]{Lemma}
\newtheorem{corollary}[theorem]{Corollary}
\newtheorem{inputtheorem}[theorem]{Input theorem}
\theoremstyle{definition}
\newtheorem{definition}[theorem]{Definition}
\theoremstyle{remark}
\newtheorem{remark}[theorem]{Remark}


\title{Microscopic--Macroscopic Independence and\\
Poisson Peeling for Luce Permutations}
\author{Christopher D. Long\\\texttt{galizur@gmail.com}}
\date{20 October 2026}

\begin{document}
\maketitle
\begin{abstract}
We study the dependence between different cycle scales of Luce permutations.
For rates converging uniformly to a positive continuous profile, every fixed
vector of short-cycle counts is asymptotically independent of the entire
ranked macroscopic $\mathrm{PD}(1)$ partition, with convergence of mixed
factorial moments against bounded continuous partition observables.
The proof verifies stability under finitely many prescribed cycles on good
clock-slab data and then averages with factorial-moment weights.
For arbitrary comparable-rate arrays, a marked switching argument gives a
quantitative joint Poisson approximation for two separated cycle-length
annuli with the remaining count as a passive mark.  Its finite-depth error
bounds can be iterated uniformly over logarithmic scales; fixed-moment
transfer is justified separately from total variation.
An appendix gives an analytic cutoff law, with an explicit error,
for the number of parts of the limiting $\mathrm{PD}(1)$ partition.
The main results concern joint dependence and complement, rather than
reprove, the scalar total-cycle analytic mod-Poisson theorem.
\end{abstract}
\setcounter{tocdepth}{1}
\tableofcontents
\newpage
\raggedbottom

\section{Model, scope, and main results}\label{sec:introduction}

Chatterjee--Diaconis--Kim~\cite[Introduction]{ChatterjeeDiaconisKim2023}
and Borga--Chatterjee--Diaconis~\cite[Section~5.2]{BorgaChatterjeeDiaconis2025}
ask for the cycle structure of Luce permutations, including joint cycle
counts.  Separate Poisson laws for short cycles and a Poisson--Dirichlet
law for large cycles do not specify their dependence.  The principal
result here supplies this missing joint statement for regular positive
profiles: every fixed short-cycle vector is asymptotically independent of
the entire ranked macroscopic partition.  This answers a part of the
broader joint-cycle question; it does not claim universality for all
positive-weight arrays.

Let $X_{n,i}\sim\operatorname{Exp}(\theta_{n,i})$, $1\le i\le n$, be
independent, and set
\[
 \sigma_n(i)=1+\#\{j:X_{n,j}<X_{n,i}\}.
\]
This is the label-to-rank Luce permutation.  Its inverse is the sequential
weighted draw order and has the stationary Tsetlin-library distribution;
inversion preserves all cycle counts.  See Luce~\cite{Luce1959},
Chatterjee--Diaconis--Kim~\cite{ChatterjeeDiaconisKim2023}, and
Borga--Chatterjee--Diaconis~\cite{BorgaChatterjeeDiaconis2025}.
Write $C_{n,r}$ for the number of $r$-cycles and $L_n(i)$ for the cycle
length at label $i$.  Put
\[
 K_n=\sum_{r=1}^nC_{n,r},\qquad
 \mathbf P_n=(L_{n,1}/n,L_{n,2}/n,\ldots),
\]
where $L_{n,1}\ge L_{n,2}\ge\cdots$ are the ranked cycle lengths.
The partition space
\[
 \mathcal S^\downarrow=
 \{(x_j)_{j\ge1}:x_1\ge x_2\ge\cdots\ge0,\ \sum_jx_j=1\}
\]
is equipped with the $\ell^1$ metric.

We distinguish two hypotheses.  The quantitative peeling results assume
only a fixed bound on the largest-to-smallest rate ratio; after a common
rescaling, write
\begin{equation}\label{eq:comparable}
 1\le\theta_{n,i}\le\varkappa<\infty.
\end{equation}
For the short-cycle identification and joint microscopic--macroscopic
limit, we assume
\begin{equation}\label{eq:profile}
 \theta_{n,i}=a_n\{f(i/n)+\varepsilon_{n,i}\},\qquad
 \max_i|\varepsilon_{n,i}|\longrightarrow0,\qquad
 f\in C[0,1],\quad 0<\min f\le\max f<\infty.
\end{equation}
The factor $a_n>0$ rescales time only.  In particular,
\eqref{eq:profile} implies \eqref{eq:comparable} with a fixed enlarged
constant for all sufficiently large $n$.

Define
\[
 \mathcal A_f(t)=\int_0^1(1-e^{-f(x)t})\,dx,\qquad
 \mathcal B_f(t)=\mathcal A_f'(t),\qquad \tau_f=\mathcal A_f^{-1},
\]
and the microscopic kernel and fixed-cycle means
\begin{align}
 \rho_f(x,y)&=\frac{f(x)e^{-f(x)\tau_f(y)}}{\mathcal B_f(\tau_f(y))},
       \qquad 0<y<1,\label{eq:rho}\\
 \lambda_r&=\frac1r\int_{[0,1]^r}
     \rho_f(x_1,x_2)\cdots\rho_f(x_r,x_1)\,dx_1\cdots dx_r.
     \label{eq:lambda}
\end{align}
The cyclic integral is interpreted by increasing bulk truncations.
Its finiteness and the fixed-cycle Poisson laws are companion-paper inputs,
not new assertions of this paper.
For falling factorials, $(x)_j=x(x-1)\cdots(x-j+1)$ and $(x)_0=1$.
We use $d_{\mathrm{TV}}$ for total variation (half the $\ell^1$ distance
on countable spaces).  For a nonnegative integer variable with the required
moments, $\kappa_q^{\mathrm F}$ denotes the factorial cumulant: the
$q$th coefficient derivative of the formal series
$\log\E(1+w)^Y$, or the corresponding analytic derivative when it exists.
Only fixed orders are used in the peeling estimates.

\begin{theorem}[Short cycles and the macroscopic partition]
\label{prop:micro-macro}
Under \eqref{eq:profile}, for every fixed $B$,
\[
 (C_{n,1},\ldots,C_{n,B},\mathbf P_n)
 \Longrightarrow (Z_1,\ldots,Z_B,\mathbf P)
\]
in $\mathbb N_0^B\times\mathcal S^\downarrow$, where
$Z_r\sim\Pois(\lambda_r)$ are independent, $\mathbf P\sim\PD(1)$,
and the Poisson vector and partition are independent.
For every fixed nonnegative integer vector $(k_1,\ldots,k_B)$ and
bounded continuous $F:\mathcal S^\downarrow\to\mathbb R$,
\begin{equation}\label{eq:micro-macro-factorization}
 \E\left[F(\mathbf P_n)\prod_{r\le B}(C_{n,r})_{k_r}\right]
 \longrightarrow \E F(\mathbf P)\prod_{r\le B}\lambda_r^{k_r}.
\end{equation}
\end{theorem}

The proof in Section~\ref{sec:micro-macro} uses a uniform conditional
macroscopic statement on good slab data.  Merely combining the two
marginal convergence theorems would not prove independence.  A consequence
is uniform mixed-moment convergence for macroscopic threshold counts on
compact threshold intervals.

Our second result, Theorem~\ref{lem:annular}, is a finite quantitative
statement under \eqref{eq:comparable}.  For nested cutoffs
$b'<b<U<U'$, it compares the \emph{joint} law of the lower-annulus count,
the upper-annulus count, and the retained count with two independent
Poisson variables independent of a copy of the retained count.  The
retained marginal is preserved exactly.  Section~\ref{sec:peeling}
then gives stage-uniform iterated-logarithm estimates and explains the
separate moment assumptions needed for unbounded tests.

Appendix~\ref{sec:PD-cumulants} treats a separate property of the limiting
$\PD(1)$ partition.  It proves an analytic threshold-count limit with an
explicit cutoff error by an elementary recurrence.  This is neither a
finite-$n$ Luce convergence rate nor another proof of the scalar
total-cycle theorem.

\paragraph{Relation to the scalar total-cycle paper.}
The profile-dependent mean shift, trace-class and Fredholm formula,
universal factorial cumulants of $K_n$, and its local analytic mod-Poisson
limit are proved in the companion manuscript~\cite{LongModPoisson2026}.
We neither reprove these results nor treat analytic mod-Poisson convergence
as an unresolved scalar problem.  That manuscript is cited for context
only: no result from it is a premise below.  Conversely, its scalar
transform alone does not imply the joint independence statement or the
passive-mark estimates proved here.  Quantitative total-cycle Poisson and
lattice approximations are treated separately in~\cite{LongTotalCycles2026}.

\section{Companion inputs and switching estimates}\label{sec:inputs}

The named companion sources in the bibliography are pinned to a fixed
repository revision.  We state the interfaces used, including the
conditional macroscopic information; unconditional $\PD(1)$ convergence
alone would not be sufficient for Section~\ref{sec:micro-macro}.

\subsection{Microscopic input}
\begin{inputtheorem}[Fixed short cycles]\label{input:micro}
Under \eqref{eq:profile}, every fixed vector $(C_{n,1},\ldots,C_{n,B})$
converges to independent Poisson variables with means \eqref{eq:lambda},
and every fixed mixed factorial moment converges to the corresponding
Poisson moment.
\end{inputtheorem}
This is the regular-profile fixed-cycle theorem and mixed factorial-moment
proposition of~\cite[Section~3]{LongMicroscopic2026}.  That paper's local
entry estimate is uniform in arbitrary distinct source labels; bulk and
separation restrictions are imposed on the target ranks only.  Its
endpoint and collision envelopes justify the globally summed moments.
We use the moment conclusion here, not a new marked local-entry argument.

\subsection{Switching and a finite endpoint estimate}

Assume \eqref{eq:comparable}.
For $i\ne j$, put
\[
 R_{ij}=\exp\{-(\theta_{n,i}-\theta_{n,j})(X_{n,j}-X_{n,i})\},
 \qquad R_{ii}=1,
\]
and define
\[
 S_{n,r}=\sum_{i:L_n(i)>r}R_{i,\sigma_n^r(i)},\qquad
 e_{n,r}=\E|S_{n,r}/n-1|,\qquad
 a_{n,r}=\frac{n}{r(n-r)}.
\]
Here $L_n(i)$ is the length of the cycle containing $i$.
Fix $p=1+1/(4\varkappa)$ when $\varkappa>1$, and take $p=2$
when $\varkappa=1$.  Then $p<\varkappa/(\varkappa-1)$ in the former
case.  Set $\beta=1-1/p$; this choice makes all constants below depend
only on the displayed parameters and on $\varkappa$.

\begin{inputtheorem}[Marked switch and selected-factor control]
\label{input:switch}
For every bounded permutation-measurable $F$,
\begin{equation}\label{eq:master-switch}
 r(n-r)\E[C_{n,r}F(\sigma_n)]
 =\E\sum_{i:L_n(i)>r}
 R_{i,\sigma_n^r(i)}
 F\bigl(\sigma_n\circ(i\ \sigma_n^r(i))\bigr).
\end{equation}
Moreover,
\begin{equation}\label{eq:selected-p}
 \sup_{n,r}\frac1n\E\sum_{i:L_n(i)>r}
 R_{i,\sigma_n^r(i)}^p<\infty.
\end{equation}
The analogous $p$-moment bound holds for an independent uniform pair of
labels.  For such a pair $(I,J)$, $\E R_{IJ}=1$.
\end{inputtheorem}

These are the marked master identity and selected-moment bounds of
\cite{LongMesoscopic2026}.  They immediately imply the following useful
pointwise version of the short-root estimate.

\begin{lemma}[Root point masses and exceptional switches]
\label{lem:root-mass}
For $1\le r<n$,
\begin{equation}\label{eq:crude-cycle}
 \E C_{n,r}\le \frac{C_\varkappa n}{r(n-r)}.
\end{equation}
Consequently, for an independent uniform label $I_n$ and any set
$E\subseteq\{1,\ldots,\lfloor n/2\rfloor\}$,
\begin{equation}\label{eq:root-set}
 \Pp\{L_n(I_n)\in E\}\le C_\varkappa |E|/n.
\end{equation}
For every set $E$ of cycle lengths,
\begin{equation}\label{eq:exception-holder}
 \frac1n\E\sum_{\substack{i:L_n(i)>r\\L_n(i)\in E}}
 R_{i,\sigma_n^r(i)}
 \le C_\varkappa\bigl(\Pp\{L_n(I_n)\in E\}\bigr)^\beta.
\end{equation}
\end{lemma}
\begin{proof}
Apply \eqref{eq:master-switch} with $F=1$ and use
\eqref{eq:selected-p} to bound $\E S_{n,r}$ by $C_\varkappa n$.
Since
\[
 \Pp\{L_n(I_n)=r\}=r\E C_{n,r}/n,
\]
\eqref{eq:crude-cycle} gives a bound $C_\varkappa/n$ for every $r\le n/2$.
Summation proves \eqref{eq:root-set}.  H\"older's inequality on the product
of clock probability and the normalized counting measure on labels proves
\eqref{eq:exception-holder}.
\end{proof}

We record the quantitative endpoint interface, including its truncation
parameter.  This avoids an appeal to an asymptotic admissibility condition
at the terminal stages of the peeling argument.
For a window $[r_-,r_+]$ with $r_+\le\delta_0 n/2$, set
\[
 \delta=2r_+/n,\qquad h=\delta^{1/4},\qquad
 \omega=\delta^{1/(4\varkappa)},\qquad M=n\sqrt\delta,\qquad
 L_\delta=1+\log(1/\delta).
\]
The factor two accommodates the two-root construction.  For an integer
terminal depth $d$, define
\begin{equation}\label{eq:endpoint-bracket}
\begin{split}
 \mathcal B_n(d;r_-,r_+)={}&
 (1-\alpha)^{d-1}+dh+d\omega+\frac{d^2+d\log n}{M}
 +\omega+\delta+\delta^{1/2}L_\delta\\
 &+h^{-1}L_\delta e^{-cM}
   +n^2h^{-1}e^{-chM}+n^{-2}.
\end{split}
\end{equation}
Here $c>0$ and $\alpha\in(0,1)$ depend only on $\varkappa$.

\begin{inputtheorem}[Finite truncation-uniform endpoint interface]
\label{input:endpoint}
There are positive constants $c_0,\delta_0$ depending only on $\varkappa$ such
that the following bounds hold for all sufficiently large $n$, provided
\[
 1\le d\le r_-/3,\qquad d\le c_0hM,\qquad \delta\le\delta_0.
\]
For every $T\ge1$, put
\[
 Y_{n,r}^{(T)}=\frac1n\sum_{i:L_n(i)>r}(R_{i,\sigma_n^r(i)}\wedge T),
 \qquad D_{n,T}=\frac1{n^2}\sum_{i,j}(R_{ij}\wedge T).
\]
Uniformly for $r\in[r_-,r_+]$,
\begin{align}
 |\E Y_{n,r}^{(T)}-\E D_{n,T}|
 &\le C_\varkappa T\mathcal B_n,
 \label{eq:endpoint-first}\\
 |\E(Y_{n,r}^{(T)})^2-\E D_{n,T}^2|
 &\le C_\varkappa T^2\mathcal B_n,
 \label{eq:endpoint-second}\\
 \Var D_{n,T}&\le 4T^2/n.
 \label{eq:endpoint-independent-var}
\end{align}
The constants are independent of $T,d,r_-,r_+$.
\end{inputtheorem}

This finite form is read from the estimates in the proof of the
free-depth endpoint proposition (\texttt{prop:endpoint-mixing},
equation \texttt{eq:free-depth-global-bound}) and its observable corollary
(\texttt{cor:quantitative-truncation-interface}) in
\cite{LongMesoscopic2026}.  These give
\eqref{eq:endpoint-first}--\eqref{eq:endpoint-second}.
Their finite estimates use $d\le r_-/3$ to leave the frozen prefixes intact
and $d\le c_0hM$ to preserve residual row and cell margins.  The conditions
$d\to\infty$, $dh\to0$, and $d\omega\to0$ in the asymptotic formulation are
used to make the displayed bound tend to zero, not to derive that bound.
For large $n$, $M\ge\sqrt{2n}$ and $hM\ge2^{3/4}n^{1/4}$, so the
finite environment estimates apply uniformly.  If the bracket is bounded
below by a fixed positive constant, the first two estimates also follow
from the trivial bounds $0\le Y_{n,r}^{(T)},D_{n,T}\le T$.
The last inequality follows
by replacing one clock: at most $2n$ summands change, each by at most $T$.
No untruncated selected second moment is required.

\begin{lemma}[Finite switching-intensity bound]\label{lem:finite-intensity}
Under the numerical conditions of Input theorem~\ref{input:endpoint},
\begin{equation}\label{eq:finite-intensity}
 \sup_{r_-\le r\le r_+} e_{n,r}
 \le C_\varkappa\bigl(\mathcal B_n(d;r_-,r_+)+n^{-1}\bigr)^\eta,
 \qquad \eta=\frac{p-1}{2p}>0.
\end{equation}
\end{lemma}
\begin{proof}
Write $\mu_{n,T}=\E D_{n,T}\le1$.  The selected and independent-pair
$p$-moment bounds give
\[
 1-\mu_{n,T}\le C_\varkappa T^{1-p},\qquad
 \E(S_{n,r}/n-Y_{n,r}^{(T)})\le C_\varkappa T^{1-p}.
\]
Using \eqref{eq:endpoint-first}--\eqref{eq:endpoint-independent-var},
\[
 \E|Y_{n,r}^{(T)}-\mu_{n,T}|^2
 \le C_\varkappa T^2(\mathcal B_n+n^{-1}).
\]
Thus, with $u=\mathcal B_n+n^{-1}$,
\[
 e_{n,r}\le C_\varkappa\{T\sqrt u+T^{1-p}\}.
\]
For $u\le1$, choose $T=u^{-1/(2p)}$.  For $u>1$, use the uniform
first-moment bound on $S_{n,r}/n$.  This proves the claim.
\end{proof}

\subsection{Fixed-order cycle domination}

\begin{inputtheorem}[Fixed-order factorial domination]
\label{input:factorial-dom}
For every fixed $j\ge1$, there is $C_{\varkappa,j}<\infty$ such that
\begin{equation}\label{eq:factorial-dom}
 \E\cC_n(r_1,\ldots,r_j)
 \le\frac{C_{\varkappa,j}}{r_1\cdots r_j}
\end{equation}
whenever $r_1+\cdots+r_j\le n/2$, where
$\cC_n(r_1,\ldots,r_j)$ counts ordered tuples of distinct cycles of the
specified lengths.
\end{inputtheorem}

This is Corollary \emph{Fixed-order factorial domination}
(\texttt{cor:factorial-domination}) of~\cite{LongOrderEdgeworth2026},
following its Exact fixed-order multi-clock identity
(\texttt{thm:multiclock}), in the range where
the total prescribed length is at most $n/2$.  Only this finite-order
probabilistic module is imported; no arithmetic order theorem or
Edgeworth expansion is used.


\subsection{Macroscopic input and its conditional scope}
\begin{inputtheorem}[Macroscopic partition and regular-history interfaces]
\label{input:PD}
Under \eqref{eq:comparable}, $\mathbf P_n\Rightarrow\PD(1)$ in $\ell^1$,
and the size-biased residual deletion sequence converges to $\GEM(1)$.
Under \eqref{eq:profile}, we also use the following conditional interfaces
of the same proof.  Normalize clocks and rates by
\[
 T_{n,i}=(\max_j\theta_{n,j})X_{n,i},\qquad
 r_{n,i}=\theta_{n,i}/\max_j\theta_{n,j}\in[\varkappa^{-1},1].
\]
For a fixed finite master partition of this normalized clock time, the
complete slab data record the label sets $H_a$, the rank sets $Q_a$, and the ordered
\emph{unlabeled} clock values.  On good slab events of probability tending
to one, each fixed row and cell population is bounded below by
$c_{\mathcal J}n$.
After conditioning only on compatible revealed edges and deleting complete
cycles, the residual matchings are independent permanental matchings.
The uniform rooted-closure and deletion estimates apply to regular residual
histories with remaining mass at least a fixed positive fraction and
relative row/cell error tending to zero.  For a partition endpoint $U$,
write $p_n(D(U))=n^{-1}\#\{i:T_{n,i}>U\}$ and
\[
 \Delta_n(U,\delta,h)
 =\delta+2h+\chi_n+\sqrt{p_n(D(U))},
\]
where $h$ is the mesh of the finite clock cells.  More precisely, if
$S<U$ are fixed clock cutoffs, the interior root-CDF error is at most
\begin{equation}\label{eq:imported-root-error}
 C_{\varkappa,S,\xi,\zeta}\Delta_n(U,\delta,h)
 +C_\xi p_n(D(S))+\pi_n+C_{\mathcal J,\xi,\zeta}/n,
\end{equation}
where $\xi$ bounds the residual fraction below, $\zeta$ is the interior
cutoff, and $\chi_n,\pi_n\to0$ for the fixed partition.  On the same good
slab events, for a deterministic $\varepsilon_n\to0$ and every partition
endpoint $U$,
\begin{equation}\label{eq:imported-tail-bound}
 p_n(D(U))\le e^{-U/\varkappa}+\varepsilon_n^2.
\end{equation}
The deletion lemma is a \emph{joint-event} bound: the probability that an
interior root cycle is deleted and the resulting history is not regular
at inner cutoff $S$ is at most $\pi_n$.  The new regularity error is
bounded by
\begin{equation}\label{eq:imported-deletion-error}
 C_{\varkappa,S,\xi,\zeta}\Delta_n(U,\delta,h)
 +C_{\mathcal J,\xi,\zeta}/n.
\end{equation}
These bounds are uniform over regular starting histories.  A uniform
finite-dimensional deletion induction from a bounded prescribed deletion
is derived below; it is not inferred from unconditional $\GEM(1)$
convergence.
\end{inputtheorem}

The unconditional conclusions are Theorem~1.1 and Corollary~8.1 of
\cite{LongPD2026}; the latter covers comparable arrays without an empirical
rate-law limit.  The conditional inputs are its residual-factorization
lemma, uniform residual-root law, and deletion-stability lemma.  The
explicit bounds above are its equations \texttt{eq:root-cdf} and
\texttt{eq:deletion-error}, together with the joint-event assertion
\texttt{eq:deletion-joint}.  Under \eqref{eq:profile}, the normalized
empirical rate law converges to the law $\nu$ of $f(U)/\max f$, with $U$
uniform on $[0,1]$.  This law is supported on $[\varkappa^{-1},1]$.  The
common-component argument of Corollary~8.1 gives
$q(\nu)\le\sqrt{1-\varkappa^{-1}}<1$, and the channel density satisfies
$B_\nu(S)=\int r e^{-rS}\nu(\dd r)\ge\varkappa^{-1}e^{-S}$.
Thus the conditional estimates apply with
fixed finite constants.  The tail bound
\eqref{eq:imported-tail-bound} follows from $r_{n,i}\ge\varkappa^{-1}$ and
concentration of the empirical clock tails at the finitely many endpoints;
the good slab event is intersected with these events.  No uniformity over
varying profile functions is needed.
The regularity tested there compares, on the rows below the chosen outer
cutoff,
\[
 \frac{|H_a\cap R|}{v_R|H_a|},\qquad
 \frac{|Q_a\cap H_b\cap R|}{v_R|Q_a\cap H_b|},
 \qquad v_R=|R|/n,
\]
with one, uniformly in the target category $b$.  The histories contain no
additional information about unexposed edges.  The source labels
\texttt{lem:residual-factorization},
\texttt{prop:uniform-root-}\allowbreak\texttt{fraction},
\texttt{lem:deletion-stability} identify the conditional modules.
The cutoff ordering follows the proof of
\texttt{prop:residual-uniforms}, but its unconditional statement is not used as a substitute
for the displayed conditional estimates.
Section~\ref{sec:micro-macro} verifies these conditions for every compatible
collection of finitely many prescribed complete cycles before invoking
their uniformity.  No bound on the probability of a bad slab event
\emph{conditional on a rare prescribed cycle} is assumed.
\section{Microscopic--macroscopic independence}\label{sec:micro-macro}

Throughout this section assume \eqref{eq:profile}.  We first establish the
finite-deletion consequence of Input~\ref{input:PD} that is needed for
weighted averaging.  All prescribed cycles in this
section have bounded total size, fixed before $n\to\infty$.

\subsection{Uniformity after finitely many prescribed cycles}

\begin{lemma}[Finite-deletion law on good slab data]
\label{lem:finite-deletion}
Fix $v<\infty$ and a bounded Lipschitz function $F$ of the ranked partition
in the $\ell^1$ metric.  Let $c_F=\E F(\mathbf P)$ for
$\mathbf P\sim\PD(1)$.  For every $\eta>0$ there are a fixed finite clock
partition, its complete slab-data sigma-field $\mathscr S_n$, and
$\mathscr S_n$-measurable events $\mathcal G_n$ with
$\Pp(\mathcal G_n^c)\to0$ such that
\begin{equation}\label{eq:slab-uniform}
 \limsup_{n\to\infty}
 \mathop{\operatorname{ess\,sup}}_{\mathscr S_n\in\mathcal G_n}
 \sup_{\substack{\Gamma\ \mathrm{compatible}\\|V(\Gamma)|\le v}}
 \left|\E[F(\mathbf P_n)\mid\mathscr S_n,\Gamma]-c_F\right|
 \le\eta.
\end{equation}
Here $\Gamma$ is a prescribed collection of disjoint complete cycles and
conditioning on $\Gamma$ means conditioning on their specified edges.
Only compatible collections of positive conditional probability are used.
\end{lemma}
\begin{proof}
Use the complete slab data and regular-history interfaces of
Input~\ref{input:PD}: the label sets $H_a$,
the rank intervals $Q_a$, and the ordered \emph{unlabeled} clock values
within each slab.  Given these data, the within-slab matchings are
independent permanental matchings.  Conditioning on compatible specified
edges removes their sources and targets and leaves the same permanental
laws on the residual minors.  This is the residual-factorization lemma of
that paper, and does not condition on the fully labeled clocks.

Write $m_a=|H_a|$, $N_{ab}=|Q_a\cap H_b|$, and let $R$ be the vertex set
remaining after deletion of $\Gamma$.  Put $r=|R|/n$,
$m_a^R=|H_a\cap R|$, and $N_{ab}^R=|Q_a\cap H_b\cap R|$.
For each deleted complete cycle $C$,
\[
 |C\cap H_a|=|C\cap Q_a|,
\]
because $\sigma_n$ maps the first set bijectively onto the second.
Thus the residual rows remain square.  For the fixed partition, the good
slab event in \cite{LongPD2026} has every row and cell population at least
$c_{\mathcal J}n$.  Since deletion changes each population by at most $v$,
\begin{equation}\label{eq:finite-deletion-regularity}
 \max_a\left|\frac{m_a^R}{r m_a}-1\right|
 +\max_{a,b}\left|\frac{N_{ab}^R}{r N_{ab}}-1\right|
 \le C_{\mathcal J}v/n.
\end{equation}
This verifies the companion's regular-history hypothesis uniformly over
$\Gamma$.

Fix the number $k$ of further uniform residual-root deletions, an interior
cutoff $\zeta\in(0,1/4)$, and a desired finite-dimensional error.
For all large $n$, the initial residual size is at least $n/2$; on the
first $k$ interior deletions it stays at least
$\xi_*n$, where $\xi_*:=\zeta^k/2$.
Choose clock cutoffs $T_0<\cdots<T_k$ from inner to outer, and then a
single master partition with mesh $h$ containing those cutoffs.
At stage $j+1$, use outer cutoff $T_{k-j}$ and inner cutoff $T_{k-j-1}$.
At each inner cutoff $T_i$, let $R_i$ and $B_i$ bound respectively the
principal root-law and deletion coefficients, with residual fraction at
least $\xi_*$, and set
\[
 A_i=\prod_{\ell=0}^i(1+R_\ell+B_\ell).
\]
All these constants are fixed before $T_{i+1}$ is chosen.  A tail error
introduced at level $j$ enters a later regularity bound with coefficient
$B_i\cdots B_j$, or a later root-law bound with coefficient
$R_iB_{i+1}\cdots B_j$, for some $i\le j$; both are at most $A_j$.
Choose $T_0$ to control the direct root tail, and each $T_{i+1}$ so that
$A_i e^{-T_{i+1}/(2\varkappa)}$ is smaller than its allocated error budget.
Only after this finite hierarchy is fixed is $h$ chosen so that
$2h\sum_{i<k}A_i$ is small.  Finally take $n$ large enough to control the
fixed-partition errors and their finitely many amplification factors.

Here is the explicit perturbation from a prescribed deletion.  Put
$d_0=C_{\mathcal J}v/n$ and, for $0\le j<k$, define
\begin{equation}\label{eq:finite-deletion-recursion}
 d_{j+1}=B_{k-j-1}
 (d_j+2h+\chi_n+\varepsilon_n+e^{-T_{k-j}/(2\varkappa)})
 +C_{\mathcal J,k,\xi_*,\zeta}/n,
\end{equation}
where $B_i$ bounds the deletion coefficient with inner cutoff $T_i$.
All sequences are replaced by the maxima over the finitely many cutoff
levels.  Relative to the same recursion starting from zero, the only
additional error is
\[
 d_j-d_j^{(0)}
 =\frac{C_{\mathcal J}v}{n}
   \prod_{i=0}^{j-1}B_{k-i-1}=o_n(1),
\]
uniformly in the compatible prescribed cycles, since all constants are
fixed before $n\to\infty$.  The hierarchy makes $d_j<1/4$ throughout.
Equations~\eqref{eq:imported-root-error}--\eqref{eq:imported-deletion-error}
then apply at every stage with these deterministic regularity bounds.

To pass from these conditional bounds to a joint law, first use bounded
Lipschitz one-variable tests supported in $[\zeta,1-\zeta]$ and vanishing
at the endpoints.  Conditional Stieltjes integration bounds their
expectation errors by the interior CDF error times the integral of the
absolute derivative.  Multiply by the previously exposed tests and
iterate.  The joint-event deletion bound charges failures of the next
regular history directly; it never conditions on the exact new cycle
length.  The total failure probability is at most the sum of the $k$
uniform failure bounds.  Evaluating the CDF at $2\zeta$ and $1-2\zeta$
bounds the omitted endpoint probabilities by $4k\zeta$ plus these errors.
Letting the chosen finite-dimensional error and then $\zeta$ tend to zero
proves the product-uniform comparison, uniformly over the good initial
slab data and the prescribed cycles.  Products of Lipschitz functions
and a finite uniform approximation on $[0,1]^k$ give the comparison for
any fixed continuous test of the first $k$ residual fractions.

In particular, for any prescribed $\eta_1>0$ the hierarchy can be chosen
so that the residual mass $R_k^{\mathrm{res}}$, normalized by the
initial residual size $|R|$, satisfies
\[
 \E[R_k^{\mathrm{res}}\mid\mathscr S_n,\Gamma]
 \le 2^{-k}+\eta_1+o_n(1)
\]
uniformly: apply the finite-dimensional comparison to
$\prod_{j=1}^k(1-u_j)$.
A bounded Lipschitz test on the ranked partition can therefore be reduced
to its first $k$ size-biased pieces at cost
$C_F(2^{-k}+\eta_1+o_n(1))$; decreasing rearrangement is a contraction in
$\ell^1$.  Formally, extend $F$ to the ambient $\ell^1$ space by a bounded
Lipschitz extension before truncating the sequence.

Finally, restoring the deleted pieces and changing normalization from
$|R|$ to $n$ changes the ranked sequence by at most $2v/n$ in $\ell^1$.
Choose $k$ large and then the finite-dimensional errors small enough for
the prescribed $\eta$.  This proves \eqref{eq:slab-uniform}.
The regular-profile hypothesis ensures convergence of the marked empirical
rate measures, as required for this fixed-partition construction.
\end{proof}

\begin{proof}[Proof of Theorem~\ref{prop:micro-macro}]
Set
\[
 W_n=\prod_{r\le B}(C_{n,r})_{k_r},\qquad
 v=\sum_{r\le B}rk_r.
\]
Expand $W_n$ as a sum of indicators over ordered collections $\Gamma$ of
the prescribed disjoint cycles.  The total number of exposed vertices is
$v$, independent of $n$.  By Input theorem~\ref{input:micro}, $\E W_n$ converges to
$\prod_r\lambda_r^{k_r}$.  Also $\sup_n\E W_n^2<\infty$: the square is a fixed polynomial in
falling factorials of the finitely many short-cycle counts, whose mixed
factorial moments converge by Input~\ref{input:micro}.  This part of the
proof does not require fixed-order domination for growing cycle lengths.

First let $F$ be bounded Lipschitz.  Apply
Lemma~\ref{lem:finite-deletion} with error $\eta$.
On $\mathcal G_n$, condition first on the complete slab data and then on
one prescribed collection.  Sum the resulting inequalities with their
conditional occurrence probabilities.  On the bad slab event use
Cauchy--Schwarz.  This gives
\begin{equation}\label{eq:weighted-good-slab}
\begin{split}
 |\E[F(\mathbf P_n)W_n]-c_F\E W_n|
 \le{}&(\eta+o_n(1))\E W_n\\
 &+2\|F\|_\infty(\E W_n^2)^{1/2}
                         \Pp(\mathcal G_n^c)^{1/2}.
\end{split}
\end{equation}
Let $n\to\infty$ and then $\eta\downarrow0$.  This proves the mixed
factorial-moment factorization for bounded Lipschitz tests.  It never
requires a uniform bound on $\Pp(\mathcal G_n^c\mid\Gamma)$.

For completeness, this factorization determines joint weak convergence.
Both marginals are tight.  Along any joint subsequential limit, uniform
integrability transfers the mixed factorial moments; conditioning by a
nonnegative bounded Lipschitz test of the partition leaves the same
factorial moments as the independent Poisson vector.  The Poisson vector
is moment-determinate.  For example, its factorial moments give a finite
exponential moment by the monotone expansion of
$\prod_r(1+t_r)^{Z_r}$ for $t_r\ge0$.
Thus every subsequential limit is the stated product law.
Finally, joint weak convergence and the uniform integrability of $W_n$
extend \eqref{eq:micro-macro-factorization} to bounded continuous $F$.
\end{proof}

\begin{corollary}[Uniform threshold moments]\label{cor:threshold-joint}
Fix $B$, integers $k_1,\ldots,k_B,k\ge0$, and
$[a,b]\subset(0,1)$.  Define
\[
 N_{n,t}=\sum_{r>tn}C_{n,r},\qquad
 N_t=\#\{j:P_j>t\},\qquad\mathbf P\sim\PD(1).
\]
Then, uniformly for $t\in[a,b]$,
\begin{equation}\label{eq:threshold-moments}
 \E\left[\prod_{r\le B}(C_{n,r})_{k_r}(N_{n,t})_k\right]
 \longrightarrow
 \left(\prod_{r\le B}\lambda_r^{k_r}\right)\E(N_t)_k.
\end{equation}
Consequently the sum $\sum_{r\le B}C_{n,r}+N_{n,t}$ has the same limiting
fixed factorial moments, uniformly on $[a,b]$, as $P_B+N_t$, where
$P_B\sim\Pois(\sum_{r\le B}\lambda_r)$ is independent of $N_t$.
\end{corollary}
\begin{proof}
For a fixed $t>0$, a $\PD(1)$ partition has no part exactly equal to $t$,
and the threshold-count map is continuous at every such partition in
$\ell^1$.  Moreover, $N_{n,t}\le1/a$ for $t\in[a,b]$.
Use Theorem~\ref{prop:micro-macro} and uniform integrability of the short-cycle
factor to obtain pointwise convergence in \eqref{eq:threshold-moments}.

For fixed short-cycle factorial factor $W_n\ge0$, the functions
$g_n(t)=\E[W_n(N_{n,t})_k]$ are nonincreasing.  Their limit is continuous
on $[a,b]$: the PD factorial moment measure, recalled in
Lemma~\ref{lem:PD-factorial-measure}, gives continuity by dominated convergence
on $[a,1]^k$.  A finite-grid monotonicity sandwich therefore upgrades
pointwise to uniform convergence.  Expand the factorial moments of the
sum by the binomial identity for falling factorials.  The same conclusion
holds simultaneously for any fixed finite set of $B$'s.
\end{proof}

\section{Annular passive-Poisson peeling}\label{sec:peeling}

For integers $b<U$, define
\[
 D_n(b,U)=\sum_{r\le b}C_{n,r}+\sum_{r>U}C_{n,r},\qquad
 H(a,b)=\sum_{a<r\le b}\frac1r.
\]
All moment orders in this section are fixed before the cutoffs and error
exponents are chosen.  The quantity $e_{n,r}=\E|S_{n,r}/n-1|$ is an
absolute switching-intensity error.  A bound only on
$|\E(S_{n,r}/n)-1|$, sufficient for scalar real-fugacity growth, would not
control the passive-mark test functions used here.

\subsection{Marked Stein comparison and annular errors}

\begin{lemma}[Passive-mark Stein transfer]\label{lem:passive-stein}
Let $X$ be a nonnegative integer-valued variable with finite mean, let $D$
be a countable-valued mark, and let $\lambda\ge0$.  Suppose that for every
bounded $g$ with $\|g\|_\infty\le1$,
\[
 \left|\E[Xg(X-1,D)]-\lambda\E g(X,D)\right|\le\eta,
\]
where the term with $X=0$ is defined as zero.  Then
\[
 \TV\bigl(\cL(X,D),\Pois(\lambda)\otimes\cL(D)\bigr)\le\eta.
\]
\end{lemma}
\begin{proof}
Use the immigration--death generator on the active count while keeping the
mark fixed.  For an event $B$, center its indicator by its Poisson average
at the same mark.  The semigroup solution $h_B$ has
$|h_B(x+1,d)-h_B(x,d)|\le1$: couple chains differing by one initial
particle, which survives until time $t$ with probability $e^{-t}$, and
integrate over $t\ge0$.  Testing the assumed Campbell discrepancy with
$g(x,d)=h_B(x+1,d)-h_B(x,d)$ bounds the difference of event probabilities
by $\eta$.  The centering integrates to
$(\Pois(\lambda)\otimes\cL(D))(B)$, so the mark retains its original
marginal.  This is the marked Poisson Stein argument; see
\cite{BarbourHolstJanson1992,LongTotalCycles2026}.
\end{proof}

\Needspace{22\baselineskip}
\begin{theorem}[Joint annular peeling]\label{lem:annular}
Assume \eqref{eq:comparable}.  Let positive integer cutoffs satisfy
\[
 b'<b<U<U'<n/4,\qquad U+b<n/2.
\]
Put
\[
 J^- =\sum_{b'<r\le b}C_{n,r},\qquad
 J^+ =\sum_{U<r\le U'}C_{n,r},\qquad
 \Delta(I)=\sum_{r\in I}a_{n,r}e_{n,r}.
\]
Let
\[
 \lambda^- =\sum_{b'<r\le b}a_{n,r},\qquad
 \lambda^+ =\sum_{U<r\le U'}a_{n,r}.
\]
There are independent $Z^-\sim\Pois(\lambda^-)$,
$Z^+\sim\Pois(\lambda^+)$ and
$\widetilde D_n(b',U')\overset d=D_n(b',U')$ such that
\begin{equation}\label{eq:annular-error}
\begin{split}
 &\TV\bigl(\cL(J^-,J^+,D_n(b',U')),
       \cL(Z^-,Z^+,\widetilde D_n(b',U'))\bigr)\\
 &\quad\le C_\varkappa\left\{
 \Delta((b',b])+\Delta((U,U'])
 +H(b',b)(b/n)^\beta+H(U,U')(U'/n)^\beta\right\}.
\end{split}
\end{equation}
\end{theorem}
\begin{proof}
For the lower annulus use the countable-valued vector
$M^-=(J^+,D_n(b',U'))$ as passive mark.
A split at distance $r\in(b',b]$ increases $J^-$ by one and leaves
$M^-$ unchanged whenever its source-cycle length is outside
\[
 E_r^-=(r,r+b]\cup(U,U+r]\cup(U',U'+r].
\]
Indeed, outside the first interval both the source and residual cycles are
longer than $b$; the other two exclusions keep them on the same sides
of both $U$ and $U'$.  The $U'$ exclusion is needed because we retain
the upper-annulus count and the final remainder separately.
The set $E_r^-$ has at most $b+2r\le3b$ integer points and is contained in
$\{1,\ldots,\lfloor n/2\rfloor\}$: its largest possible endpoint is
$U'+r<U'+b<2U'<n/2$.
For instance, $b'=2,b=5,U=12,U'=20,r=3$ and source length $22$
show why the last interval cannot be omitted: the split creates the
upper-annulus cycle of length $19$ and removes a long retained cycle.
This witness is admissible for any $n>80$.

Apply \eqref{eq:master-switch} with the test
$g(J^-(\pi)-1,M^-(\pi))$ on permutations with $J^-(\pi)>0$, and zero
otherwise.  The discrepancy from the unchanged test is at most two on
exceptional splits.  Thus, for $\|g\|_\infty\le1$,
\begin{align*}
 &\left|\E[C_{n,r}g(J^--1,M^-)]-a_{n,r}\E g(J^-,M^-)\right|\\
 &\quad\le a_{n,r}e_{n,r}
 +\frac{2a_{n,r}}n
    \E\sum_{\substack{i:L_n(i)>r\\L_n(i)\in E_r^-}}
                  R_{i,\sigma_n^r(i)}\\
 &\quad\le a_{n,r}e_{n,r}+C_\varkappa a_{n,r}(b/n)^\beta,
\end{align*}
where the last step is Lemma~\ref{lem:root-mass}.
Summing over $r$ and applying Lemma~\ref{lem:passive-stein} gives a joint
approximation of $(J^-,M^-)$ by
$\Pois(\lambda^-)\otimes\cL(M^-)$ with the first and third errors in
\eqref{eq:annular-error}; here $a_{n,r}\le2/r$.

Next apply the same argument to $J^+$ with passive mark $D_n(b',U')$.
For $r\in(U,U']$, the split preserves this mark and adds one to $J^+$
whenever its source length exceeds $r+U'$.
The sufficient exceptional set $(r,r+U']$ has at most $U'$ points and lies
below $n/2$, since $r\le U'$ and $2U'<n/2$.  It is deliberately
conservative: no assertion is needed about splits inside that set.  The same argument gives the second and fourth
errors, with parameter $\lambda^+$.  Tensoring this second approximation
with the independent $Z^-$ and applying the triangle inequality proves
\eqref{eq:annular-error}.  Applying the summation map also gives the
corresponding comparison of $D_n(b,U)$ with
$Z^-+Z^++\widetilde D_n(b',U')$.
\end{proof}

\subsection{Moment bounds and factorial-cumulant transfer}

\begin{lemma}[Harmonic moment bound]\label{lem:harmonic-moments}
Fix $k\ge1$ and set $W=\lfloor n/(2k)\rfloor$.  For any set $I$ of
cycle lengths, let $V_I=\sum_{r\in I}C_{n,r}$ and
$h_I=\sum_{r\in I,\ r\le W}1/r$.  Then, for all sufficiently large $n$,
\begin{equation}\label{eq:harmonic-moments}
 \E V_I^k\le C_{\varkappa,k}(1+h_I)^k.
\end{equation}
In particular, $\E K_n^k\le C_{\varkappa,k}(1+\log n)^k$.
\end{lemma}
\begin{proof}
Let $V=\sum_{r\in I,\ r\le W}C_{n,r}$.  For $1\le j\le k$, expansion
into ordered distinct-cycle tuples and Input theorem~\ref{input:factorial-dom} give
\[
 \E(V)_j\le C_{\varkappa,j}h_I^j,
\]
because each tuple has total length at most $jW\le n/2$.
The expansion of $V^k$ in falling factorials proves the corresponding
ordinary-moment bound.  There are fewer than $2k$ cycles longer than $W$,
since each such cycle has length greater than $n/(2k)$.  Adding this
bounded count proves \eqref{eq:harmonic-moments}.
\end{proof}

\begin{lemma}[Factorial-cumulant transfer]\label{lem:cumulant-transfer}
Fix $q\ge2$.  If $X,Y$ are nonnegative integer-valued variables with
\[
 \TV(\cL(X),\cL(Y))\le\eta,\qquad
 \E X^{2q}+\E Y^{2q}\le C_q H^{2q},\qquad H\ge1,
\]
then
\begin{equation}\label{eq:cumulant-transfer}
 \left|\kappa_q^{\mathrm F}(X)-\kappa_q^{\mathrm F}(Y)\right|
 \le C_q'H^q\eta^{1/2}.
\end{equation}
\end{lemma}
\begin{proof}
Under a maximal coupling, $\Pp(X\ne Y)\le\eta$.  For $1\le j\le q$,
Cauchy--Schwarz and the $2q$-th-moment bound imply
\[
 |\E(X)_j-\E(Y)_j|
 \le\E[(X^j+Y^j)\one_{\{X\ne Y\}}]
 \le C_{q,j}H^j\eta^{1/2}.
\]
Writing $\cP_q$ for the set of partitions of $\{1,\ldots,q\}$,
the factorial cumulant is the fixed polynomial
\[
 \kappa_q^{\mathrm F}(X)
 =\sum_{\pi\in\cP_q}
 (-1)^{|\pi|-1}(|\pi|-1)!
 \prod_{B\in\pi}\E(X)_{|B|}.
\]
Telescoping each product and using $\sum_{B\in\pi}|B|=q$ proves the claim.
The constants depend on the fixed order; no uniformity in $q$ is asserted.
\end{proof}

\subsection{Uniform iterated-logarithm estimates}

Define successive logarithms while they are positive by
\[
 \ell_1=\log n,\qquad \ell_{j+1}=\log\ell_j,
\]
and, for a fixed exponent $A$,
\[
 b_j=\lceil\ell_j^3\rceil,\qquad
 U_j=\lfloor n/\ell_j^A\rfloor,\qquad D_j=D_n(b_j,U_j).
\]

\begin{lemma}[Iterated-log peeling estimate]\label{lem:iterated-peel}
Fix $q\ge2$ and $Q>2q+4$.  One can choose $A$ sufficiently large so that,
for some $H_0$ and all sufficiently large $n$, the following assertions
hold simultaneously at every stage with $x:=\ell_{j+1}\ge H_0$:
\begin{align}
 &\TV\bigl(\cL(D_j),
    \cL(Z_j^-+Z_j^++\widetilde D_{j+1})\bigr)
      \le C_{\varkappa,Q,A}x^{-Q},
 \label{eq:iterated-tv}\\
 &\E D_j^{2q}
  +\E(Z_j^-+Z_j^++\widetilde D_{j+1})^{2q}
      \le C_{\varkappa,q,A}(1+x)^{2q},
 \label{eq:iterated-moments}
\end{align}
where $Z_j^-,Z_j^+,\widetilde D_{j+1}$ are independent,
$\widetilde D_{j+1}\overset d=D_{j+1}$, and the Poisson parameters are
\[
 \lambda_j^- =\sum_{b_{j+1}<r\le b_j}a_{n,r},\qquad
 \lambda_j^+ =\sum_{U_j<r\le U_{j+1}}a_{n,r}.
\]
Consequently,
\begin{equation}\label{eq:cumulant-step}
 \left|\kappa_q^{\mathrm F}(D_j)-\kappa_q^{\mathrm F}(D_{j+1})\right|
 \le C_{\varkappa,q,Q,A}x^{q-Q/2}.
\end{equation}
\end{lemma}
\begin{proof}
Since $\ell_j=e^x$, harmonic summation gives
\begin{align}
 H(b_{j+1},b_j)&=3(x-\log x)+O(1),\notag\\
 H(U_j,U_{j+1})&=A(x-\log x)+O(1).
 \label{eq:annular-harmonic-budget}
\end{align}
In particular, the harmonic mass of either annulus is $O_A(x)$.
For large $H_0$ and then large $n$, all cutoff inequalities in
Theorem~\ref{lem:annular} hold uniformly.  Indeed,
$b_j\le2(\log n)^3$, $U_j\ge n/(2(\log n)^A)$, and
$U_{j+1}/n\le x^{-A}$.

Choose a fixed $D>0$ and terminal depth $d=\lceil D\log x\rceil$.
On the lower annulus, $r_-\ge b_{j+1}\ge x^3$; on the upper annulus,
$r_-\ge U_j$.  Thus $d\le r_-/3$ once $H_0$ is large.
Also $d=O_D(\log n)$ and $hM\ge c n^{1/4}$, so the numerical residual
margin condition in Input theorem~\ref{input:endpoint} holds uniformly.

For clarity, retain the finite-$n$ error in the calculation.  The reset
term is at most $Cx^{-cD}$.  For the lower annulus, its largest relative
length is $O((\log n)^3/n)$; hence all nonreset terms in
\eqref{eq:endpoint-bracket}, after applying
Lemma~\ref{lem:finite-intensity}, are bounded by
$n^{-c_1}(\log n)^{C_1}$ for some $c_1>0$.
For the upper annulus, its largest relative length is $O(x^{-A})$.
The terms $dh,d\rho,\rho,\delta,\delta^{1/2}L_\delta$ are bounded by
\[
 C_{\varkappa,A,D}(1+\log x)^C x^{-a_0 A}
\]
for some $a_0>0$ after the same fixed positive power is applied.
The terms $(d^2+d\log n)/M$ and the environment failures again contribute
at most $n^{-c_1}(\log n)^{C_1}$, after changing constants.
The selected-factor exceptional terms are bounded using
$(b_j/n)^\beta$ and $x^{-A\beta}$.

Thus Theorem~\ref{lem:annular} and \eqref{eq:annular-harmonic-budget} give
\begin{equation}\label{eq:explicit-peel-budget}
\begin{split}
 \text{error in \eqref{eq:iterated-tv}}
 \le C_{\varkappa,A,D}x\bigl[&x^{-c_2D}+(1+\log x)^{C_2}x^{-a_0A}
       +x^{-A\beta}\bigr]\\
 &+C_{\varkappa,A,D}x\,n^{-c_1}(\log n)^{C_1}.
\end{split}
\end{equation}
Choose $D$ with $c_2D>Q+3$, and then $A$ with
$\min(a_0,\beta)A>Q+3$.  Since every stage has
$x\le\log\log n$, the last term is $o(x^{-Q})$ uniformly.
The spare powers absorb the logarithmic factors.  This proves
\eqref{eq:iterated-tv}; constants are independent of the number of stages.

For the moments use Lemma~\ref{lem:harmonic-moments} with $k=2q$ and
$W=\lfloor n/(4q)\rfloor$.  For the lengths contributing to $D_j$,
\[
 \sum_{r\le b_j}\frac1r+
 \sum_{U_j<r\le W}\frac1r
 \le(A+3)x+O_q(1).
\]
The same bound holds for $D_{j+1}$, since its corresponding scale is
$\log x\le x$.  Finally, by $a_{n,r}\le2/r$ and
\eqref{eq:annular-harmonic-budget},
$\lambda_j^-+\lambda_j^+=O_A(x)$.
Fixed Poisson moments are polynomials in the mean, which proves
\eqref{eq:iterated-moments}.  Apply Lemma~\ref{lem:cumulant-transfer} and
additivity under independence.  Both Poisson factorial cumulants vanish
for $q\ge2$, giving \eqref{eq:cumulant-step}.
\end{proof}


\begin{corollary}[Accumulated peeling error]\label{cor:accumulated}
Fix the parameters of Lemma~\ref{lem:iterated-peel} and $H\ge H_0$.
For all sufficiently large $n$, let $J=J(n,H)$ be the last index with
$\ell_J\ge H$.  There are independent Poisson variables with means
$\lambda_j^-,\lambda_j^+$, $1\le j<J$, independent of a variable
$\widetilde D_J\overset d=D_J$, such that
\begin{equation}\label{eq:accumulated-tv}
 \TV\left(\cL(D_1),
 \cL\left(\sum_{j<J}(Z_j^-+Z_j^+)+\widetilde D_J\right)\right)
 \le C_{\varkappa,Q,A}H^{-Q}.
\end{equation}
For the fixed factorial-cumulant order $q$ of that lemma,
\begin{equation}\label{eq:accumulated-cumulant}
 |\kappa_q^{\mathrm F}(D_1)-\kappa_q^{\mathrm F}(D_J)|
 \le C_{\varkappa,q,Q,A}H^{q-Q/2}.
\end{equation}
\end{corollary}
\begin{proof}
Iterate the scalar consequence of Theorem~\ref{lem:annular}; convolution
with a fixed probability law contracts total variation.  At every step
use independent copies in the comparison law.  The total error is bounded
by $C\sum_{j<J}\ell_{j+1}^{-Q}$.  Read this finite sequence backward.
Successive values are related by $x\mapsto e^x$ and, for $H_0$ large,
are at least twice their predecessors.  Thus the sum is at most
$CH^{-Q}$, with a constant independent of $J$.  The same argument applied
to \eqref{eq:cumulant-step}, using $q-Q/2<0$, proves the second bound.
\end{proof}

\begin{remark}[Meaning of the iterated comparison]
Equation~\eqref{eq:accumulated-tv} is a comparison of the scalar retained
count with a convolution law.  It does not assert simultaneous independence
of all actual annular counts at all stages.  The one-stage vector estimate
in Theorem~\ref{lem:annular} is stronger than its scalar consequence, but
iteration with additional retained coordinates would require preserving
those coordinates in every switching test.  Neither the number of
factorial-cumulant orders nor the number of arbitrary passive coordinates
is allowed to grow under the fixed-order statements above.
\end{remark}
\section{Scope and further questions}\label{sec:open}

The independence theorem concerns every fixed short-cycle vector and the
whole ranked partition; the quantitative peeling estimates concern
unweighted comparable-rate arrays.  The explicit cutoff error in
Theorem~\ref{thm:PD-analytic} belongs to the limiting $\PD(1)$ model.
It is not a convergence rate for $\mathbf P_n$, nor does it justify
interchanging a moving microscopic cutoff with $n\to\infty$.

A natural extension is to allow the number of observed short-cycle lengths
or the size of the prescribed deletion to grow.  This would require
quantitative microscopic moment bounds and a residual-history estimate
uniform in that growing size.  Another question is a joint total-cycle
and marked short-cycle generating-function limit under a fixed real or
complex cycle tilt.  The scalar result of~\cite{LongModPoisson2026} and
the untilted joint limit here do not, by themselves, supply such a theorem:
one must also control the tilted dependence and unbounded generating tests.

The positive-profile assumption in the independence theorem excludes
endpoint-singular families.  Extending these joint statements to such
families requires a separate boundary-process analysis; no transfer is
asserted from comparable rates alone when the rate ratio is unbounded.

\paragraph{Dependency boundary.}
The independence theorem uses the microscopic mixed factorial moments and
the conditional regular-history interfaces of the macroscopic paper.
The annular and iterated comparisons additionally use the exact marked
switch, finite two-root endpoint interface, and the fixed-order factorial
domination stated in Section~\ref{sec:inputs}.  Their quantitative scope is
recorded there.  The $\PD(1)$ cutoff argument uses the standard size-biased
representation and the cited Dickman density.  The scalar mean-shift,
Fredholm, and total-cycle cumulant arguments from the earlier manuscript
are not dependencies and are not reproduced.  The present draft has been
internally checked against these interfaces; no claim of independent
external refereeing or formal verification is made.


\appendix
\section{An analytic cutoff law for the limiting PD partition}\label{sec:PD-cumulants}

Let $\mathbf P=(P_1,P_2,\ldots)\sim\PD(1)$ and
$N_\eps=\#\{j:P_j>\eps\}$.

This appendix is independent of the Luce conditioning and switching
arguments.  Its error is a function of the cutoff in the limiting
partition, not of the finite permutation size.

\begin{theorem}[Analytic cutoff law for $\PD(1)$]\label{thm:PD-analytic}
Let $\mathbf P\sim\PD(1)$ and
$N_\varepsilon=\#\{j:P_j>\varepsilon\}$.
For every $0<\varrho<1/2$, uniformly for complex $|w|\le\varrho$,
\begin{equation}\label{eq:PD-analytic}
 \varepsilon^w\E(1+w)^{N_\varepsilon}
 =\frac{e^{-\gamma w}}{\Gamma(1+w)}
   +O_\varrho(\varepsilon^{1-2\varrho})
 \qquad(\varepsilon\downarrow0).
\end{equation}
Here $\varepsilon^w=\exp(w\log\varepsilon)$ uses the real logarithm,
and $\gamma$ is Euler's constant.
\end{theorem}

\begin{lemma}[Factorial moment measures of $\PD(1)$]
\label{lem:PD-factorial-measure}
For every $k\ge1$, the ordered $k$-point factorial moment measure of the
parts is
\begin{equation}\label{eq:PD-factorial-measure}
 \one_{\{x_1+\cdots+x_k<1\}}
 \frac{\dd x_1\cdots\dd x_k}{x_1\cdots x_k}.
\end{equation}
\end{lemma}
\begin{proof}
Successive size-biased sampling of the parts gives the $\GEM(1)$ sequence
$V_j=U_j\prod_{i<j}(1-U_i)$ with independent uniform $U_i$; see
\cite[Chapter~3]{Pitman2006}.  If $s_j=x_1+\cdots+x_j$ and $s_0=0$, the
joint density of its first $k$ terms is
$\prod_{j=1}^k(1-s_{j-1})^{-1}$ on $s_k<1$.
Conditional on a mass partition, an ordered tuple of distinct parts
$(x_1,\ldots,x_k)$ is selected with probability
$\prod_{j=1}^k x_j/(1-s_{j-1})$.
Dividing the former density by this sampling weight gives
\eqref{eq:PD-factorial-measure}.  In particular,
$\E N_\eps=\log(1/\eps)$ and there is almost surely no part at a fixed
positive threshold.
\end{proof}

\begin{proof}[Proof of Theorem~\ref{thm:PD-analytic}]
Put $t=\log(1/\eps)$ and write
\[
 \Phi(t,w)=\E(1+w)^{N_{e^{-t}}},\qquad
 \Psi(t,w)=e^{-wt}\Phi(t,w).
\]
For each finite $t$, $\Phi(t,\cdot)$ is a polynomial, since the threshold
count is bounded.  By Lemma~\ref{lem:PD-factorial-measure},
\begin{equation}\label{eq:PD-pgf-series}
 \Phi(t,w)=\sum_{k\ge0}\frac{w^k}{k!}
 \int_{\substack{x_i>e^{-t}\\x_1+\cdots+x_k<1}}
 \prod_{i=1}^k\frac{\dd x_i}{x_i}.
\end{equation}
The $k$-fold integral is at most $t^k$, so
$|\Phi(t,w)|\le e^{|w|t}$.

Differentiate the moving lower boundaries in
\eqref{eq:PD-pgf-series}.  Each boundary puts one variable equal to
$e^{-t}$; rescale the remaining variables by $1-e^{-t}$.
For $t\ge\log2$, this gives
\begin{equation}\label{eq:PD-cutoff-recurrence}
 \partial_t\Phi(t,w)=w\Phi(\log(e^t-1),w).
\end{equation}
The identity holds almost everywhere, and $\Phi$ is locally absolutely
continuous in $t$, which suffices below.  For $0\le t<\log2$, at most one
part exceeds $e^{-t}$, whence $\Phi(t,w)=1+wt$.
Thus for all $t\ge0$,
\begin{equation}\label{eq:PD-derivative-bound}
 |\partial_t\Phi(t,w)|\le |w|e^{|w|t}
\end{equation}
almost everywhere.

For $t\ge\log2$, set $d(t)=-\log(1-e^{-t})\le2e^{-t}$.  Then
$\log(e^t-1)=t-d(t)$.  By
\eqref{eq:PD-cutoff-recurrence}--\eqref{eq:PD-derivative-bound},
\begin{align}
 \partial_t\Psi(t,w)
 &=we^{-wt}\{\Phi(t-d(t),w)-\Phi(t,w)\},\notag\\
 |\partial_t\Psi(t,w)|
 &\le |w|^2d(t)e^{(|w|-\Re w)t}
 \le2\varrho^2e^{-(1-2\varrho)t},\qquad |w|\le\varrho<1/2.
 \label{eq:PD-integrable-derivative}
\end{align}
This integrable bound proves locally uniform convergence of
$\Psi(t,\cdot)$ on $|w|<1/2$, with error $O_\varrho(e^{-(1-2\varrho)t})$.
The limiting function is analytic.

It remains to identify that function.  For real $w>0$, let
$\mathcal N_{\eps,w}$ be a Poisson process on $[\eps,1]$ with intensity
$w\dd x/x$, and let $S_{\eps,w}$ be the sum of its points.
The Poisson expansion in \eqref{eq:PD-pgf-series} gives
\begin{equation}\label{eq:PD-dickman-probability}
 \eps^w\E(1+w)^{N_\eps}=\Pp(S_{\eps,w}<1).
\end{equation}
As $\eps\downarrow0$, the sums increase to the time-$w$ Dickman
subordinator, with L\'evy measure $\one_{(0,1)}(x)\dd x/x$.
Its density on $(0,1]$ is
\[
 \frac{e^{-\gamma w}}{\Gamma(w)}x^{w-1};
\]
see \cite[Theorem~1.1]{CaravennaSunZygouras2019}.
There is no atom at one, and integration gives
$e^{-\gamma w}/\Gamma(1+w)$ for the limit of
\eqref{eq:PD-dickman-probability}.
The identity theorem identifies the analytic limit on $|w|<1/2$.
Integrating \eqref{eq:PD-integrable-derivative} proves the stated rate.
\end{proof}

\begin{proposition}[Small-threshold PD cumulants]\label{prop:PD-cumulants}
For every fixed $q\ge2$,
\begin{equation}\label{eq:PD-cumulant-limit}
 \kappa_q^{\mathrm F}(N_\eps)
 \longrightarrow(-1)^{q+1}(q-1)!\zeta(q).
\end{equation}
More precisely, for every fixed $0<\varrho<1/2$ the error is
$O_{q,\varrho}(\eps^{1-2\varrho})$.
For $q=2$,
\begin{equation}\label{eq:PD-var-exact}
 \kappa_2^{\mathrm F}(N_\eps)
 =-\int_{\substack{\eps<x,y<1\\x+y>1}}\frac{\dd x\dd y}{xy},
 \qquad
 0\le\zeta(2)+\kappa_2^{\mathrm F}(N_\eps)
 \le2\log\frac1{1-\eps}.
\end{equation}
\end{proposition}
\begin{proof}
The limiting function in \eqref{eq:PD-analytic} is nonzero on
$|w|\le\varrho<1/2$.  For sufficiently small $\eps$, take logarithms with
value zero at $w=0$ and use Cauchy's formula.  This justifies
coefficientwise differentiation, with the asserted error, of
\[
 \log\E(1+w)^{N_\eps}
 =w\log(1/\eps)-\gamma w-\log\Gamma(1+w)
  +O_\varrho(\eps^{1-2\varrho}).
\]
Now
\[
 \log\Gamma(1+w)
 =-\gamma w+\sum_{q\ge2}\frac{(-1)^q\zeta(q)}q w^q.
\]
For $q=2$, the one- and two-point factorial measures give the integral in
\eqref{eq:PD-var-exact} directly.  The integral over the full region
$0<x,y<1$, $x+y>1$, equals $\zeta(2)$.
The omitted part is at most
\[
 2\int_0^\eps\frac{-\log(1-x)}x\dd x
 \le2\int_0^\eps\frac{\dd x}{1-x},
\]
which proves the final bound.
\end{proof}


\begin{thebibliography}{99}
\small\raggedright
\setlength{\itemsep}{3pt}

\bibitem{BarbourHolstJanson1992}
A.~D.~Barbour, L.~Holst, and S.~Janson,
\emph{Poisson Approximation},
Oxford Studies in Probability, vol.~2, Clarendon Press, Oxford, 1992.
\url{https://doi.org/10.1093/oso/9780198522355.001.0001}.

\bibitem{BorgaChatterjeeDiaconis2025}
J.~Borga, S.~Chatterjee, and P.~Diaconis,
\emph{Permuton and local limits for the Luce model},
arXiv:2509.07729, version~3, 29 September 2026.
\url{https://arxiv.org/abs/2509.07729v3}.

\bibitem{CaravennaSunZygouras2019}
F.~Caravenna, R.~Sun, and N.~Zygouras,
The Dickman subordinator, renewal theorems, and disordered systems,
\emph{Electron. J. Probab.} \textbf{24} (2019), paper no.~101, 1--40.
\url{https://doi.org/10.1214/19-EJP353}.

\bibitem{ChatterjeeDiaconisKim2023}
S.~Chatterjee, P.~Diaconis, and G.~B.~Kim,
Enumerative theory for the Tsetlin library,
\emph{J. Algebra} \textbf{655} (2024), 139--162.
\url{https://doi.org/10.1016/j.jalgebra.2023.08.009}.

\Needspace{5\baselineskip}
\bibitem{LongMesoscopic2026}
C.~D.~Long,
\emph{Mesoscopic cycles of comparable-rate Luce permutations: a switching proof of Conjecture 13.1},
research manuscript, 2026.
Repository \texttt{long-mathematics/luce-tsetlin-permutations},
revision \texttt{38f195dc5b30}.
\path{papers/mesoscopic-switching/mesoscopic-luce-cycles.tex}.
\href{https://github.com/long-mathematics/luce-tsetlin-permutations/blob/38f195dc5b30812ebbf084764016e857a68b9a1c/papers/mesoscopic-switching/mesoscopic-luce-cycles.tex}{Pinned source}.

\bibitem{LongMicroscopic2026}
C.~D.~Long,
\emph{Microscopic cycles and sparse entries for regular Luce profiles},
research manuscript, 2026.
Repository \texttt{long-mathematics/luce-tsetlin-permutations},
revision \texttt{38f195dc5b30}.
\path{papers/regular-profile-cycles/regular-profile-cycles.tex}.
\href{https://github.com/long-mathematics/luce-tsetlin-permutations/blob/38f195dc5b30812ebbf084764016e857a68b9a1c/papers/regular-profile-cycles/regular-profile-cycles.tex}{Pinned source}.

\bibitem{LongOrderEdgeworth2026}
C.~D.~Long,
\emph{A sharp Berry--Esseen theorem and first Edgeworth expansion for the Erd\H{o}s--Tur\'an law of comparable-rate Luce permutations},
research manuscript, 2026.
Repository \texttt{long-mathematics/luce-tsetlin-permutations},
revision \texttt{38f195dc5b30}.
\path{papers/sharp-order-edgeworth/luce-edgeworth.tex}.
\href{https://github.com/long-mathematics/luce-tsetlin-permutations/blob/38f195dc5b30812ebbf084764016e857a68b9a1c/papers/sharp-order-edgeworth/luce-edgeworth.tex}{Pinned source}.

\bibitem{LongPD2026}
C.~D.~Long,
\emph{Spectral Poisson--Dirichlet universality for Luce permutations
under arbitrary rate arrangements},
research manuscript, 20 October 2026.
Repository \texttt{long-mathematics/luce-tsetlin-permutations},
\path{papers/macroscopic-pd/spectral-poisson-dirichlet-luce.tex}.
\href{https://github.com/long-mathematics/luce-tsetlin-permutations/blob/81462993a7400f972a3d75583d447fbeccf1e7b3/papers/macroscopic-pd/spectral-poisson-dirichlet-luce.tex}{Pinned repository source}.

\bibitem{LongTotalCycles2026}
C.~D.~Long,
\emph{Poisson approximation and lattice Edgeworth expansions for total Luce cycles},
research manuscript, 2026.
Repository \texttt{long-mathematics/luce-tsetlin-permutations},
revision \texttt{38f195dc5b30}.
\path{papers/total-cycle-edgeworth/total-cycle-edgeworth.tex}.
\href{https://github.com/long-mathematics/luce-tsetlin-permutations/blob/38f195dc5b30812ebbf084764016e857a68b9a1c/papers/total-cycle-edgeworth/total-cycle-edgeworth.tex}{Pinned source}.

\bibitem{LongModPoisson2026}
C.~D.~Long,
\emph{An exact uniform component and analytic mod-Poisson convergence for Luce permutations},
research manuscript, 2026.
Repository \texttt{long-mathematics/luce-tsetlin-permutations},
revision \texttt{38f195dc5b30}.
\path{papers/mod-poisson/luce_mod_poisson.tex}.
\href{https://github.com/long-mathematics/luce-tsetlin-permutations/blob/38f195dc5b30812ebbf084764016e857a68b9a1c/papers/mod-poisson/luce_mod_poisson.tex}{Pinned source}.

\bibitem{Luce1959}
R.~D.~Luce, \emph{Individual Choice Behavior: A Theoretical Analysis},
Wiley, New York, 1959.

\bibitem{Pitman2006}
J.~Pitman, \emph{Combinatorial Stochastic Processes},
Lecture Notes in Mathematics, vol.~1875, Springer, Berlin, 2006.
\url{https://doi.org/10.1007/b11601500}.

\end{thebibliography}
\end{document}
